% Transcription — fonds Alexandre Grothendieck, shelfmark 60, batch 1 (pages 1-20).
% Established from the facsimile archives/batches/60/batch-01.pdf (PDF page k =
% archivists' page k, no cover sheet), cut from a copy of 60.pdf fetched
% through the project's relay
% (https://grothendieck-relay.onrender.com/source/60.pdf); tiles in
% archives/tiles/60/p1-p20. Per the skill transcribe-grothendieck. The folder
% is sixty-nine pages in four batches (1-20, 21-40, 41-60, 61-69), transcribed
% in parallel by four passes; this is the first, and none of the others was
% available to read.
%
% Pass: Opus 5.5 (claude-opus-5-5), 2026-09-26 — first pass, unchecked against the pages by a human.
%
% Pages skipped: 3, 5, 7, 9, 11, 13, 15, 17, 19 — the typed versos of the
% sheets (Bourbaki material: pages of rédactions n° 357 and 392, the
% « Compléments au chapitre VII d'Algèbre » n° 398, a Tribu n° 57 listing),
% several scanned upside down, with nothing in his hand: skipped without
% description by the settled rule. Page 1 is the folder's cover, otherwise
% blank, inscribed top right in his hand « Notes sur l'équation
% fonctionnelle » (underlined); it takes no \page marker and its words become
% the section heading, with a \note.
%
% Pages summarised: none. Pages transcribed: 2, 4, 6, 8, 10, 12, 14, 16, 18,
% 20 — ten rectos in ink (black, then blue-black from p. 12), written on the
% backs of the typescripts. His own pagination: none visible; the pencilled
% numbers bottom left are the archivists'. His numbered formulas form one run
% across the batch — (1) p. 10, (2) p. 12, (3) p. 14 — and are kept as
% written. Fast drafting: formulas carry, a good part of the connective prose
% does not and is left \ill{}.
%
% What the batch is about. The functional equation of L-functions in the
% geometric case, over F_p. Pp. 2-4: P_f(t) = det(1 - ft) and Z_f = 1/P_f
% as traces of Sym^nu; the product * on 1 + tA[[t]] (1 - at)*(1 - bt) =
% 1 - abt, P_{f⊗f'} = P_f * P_{f'}, and the lambda-structure: L of lambda^i f
% and sigma^i f in terms of sigma^i / lambda^i of L_f. P. 6: the contragredient,
% P_{f^v}(t) = (-t)^{eps} (det f)^{-1} P_f(t^{-1}), with the remark that
% t -> t^{-1} only makes sense on series coming from polynomials. P. 8: for a
% Q_l-sheaf (first « motif ») E over F_p of rank n and weight rho satisfying
% Weil, det f_E = eps(E) p^{n rho/2}; with E^v = E(rho), L_{E^v}(t) =
% L_E(t/p^rho). P. 10: global duality H^i_c(X, DE) = D H^{-i}(X, E),
% Rf_*(DE) = D(Rf_* E), whence (1) L_{DE}(t) = (-t)^{-chi(E)} delta(E)
% L_E(t^{-1}), chi geometric invariant, delta arithmetic invariant, delta =
% eps(E) p^{sum (-1)^i (i+rho)/2 b_i} « avec les conjectures standard ».
% P. 12: in the favourable case (b_i = b_{2n-i}) the exponent is
% ((n+rho)/2) chi(E), giving (2); with E^v = E(rho) globally, DE =
% E(rho+n)[2n] and L_{DE}(t) = L_{Rf_*E}(p^{-(rho+n)} t). P. 14: (3)
% L_E(1/(p^{n+rho} t)) = eps(E) (-p^{(n+rho)/2} t)^{chi(E)} L_E(t); with
% xi_E(t) = t^{chi(E)/2} L_E(t), xi(1/(p^{n+rho} t)) = eps'(E) xi(t),
% eps(E) = sign det f on H^n, eps'(E) = (-1)^{b_n} sign det f_{H^n}. P. 16:
% n + rho odd gives eps' = 1 (b_n even, Serre's alternating-form argument);
% n + rho even, eps' = (-1)^{eps_1} with eps_{±1} the multiplicities of ±1;
% outside the favourable case (1) relates two functions L_{DE}, L_E. Pp. 18-20:
% X smooth of pure dimension n, E pure of weight rho, f compactifiable; duality
% D Rf_!(E) = Rf_*(DE), classes cl, the term at infinity Rf_infty(E) and
% L^infty_E; the boxed equation L_E(1/(p^{n+rho} t)) = A(t) delta(E)
% (-t)^{chi(E)} L_E(t) with A(t) = L^infty_E(1/(p^{n+rho} t))^{-1} (A = 1 for X
% complete), an NB on the type (<= tau + (n-1)) and weight (<= rho + 2n) of
% L^infty, a boxed caveat « en admettant la résolution des singularités ...
% et les conjectures de Weil », and, X complete, delta(E)^2 = p^{(n+rho)chi(E)},
% « compatible avec les conjectures de Weil ».
%
% Notation fixed once for the batch: f for Frobenius acting (f_E, f_{H^i}),
% \check E for the contragredient (he writes E with a háček), D for the
% dualising functor, R f_*, R f_! with \mathbb-free R; \bar X, \bar E, \bar F
% as written — he writes F and E for the same sheaf on pp. 10, 14, 20 and both
% are kept; \varepsilon, \varepsilon', \varepsilon_i, \delta, \chi, \rho as on
% the page; \mathbb{Q}_\ell, \mathbb{F}_p; L^{\infty}_E and R f_{\infty} for his
% superscript/subscript infinity (pp. 18, 20); \operatorname{cl} for « cl ».
%
% Readings to check first: the small index diagram at the foot of p. 2; the
% exponent -eps(f) in the second boxed formula of p. 6; the words bracketed
% on p. 8 after « dépendant du motif »; the exponent (i + rho)/2 on p. 10 and
% the margin under it; « Serre » and the parenthesis on p. 16; the reading
% R f_infty (p. 18) and the first line « chi(E) = eps(Rf_!(E)) » of p. 20.

\documentclass[11pt,a4paper]{article}
% Le préambule commun à toutes les transcriptions du fonds Grothendieck.
%
% Il tient en une idée : **l'incertitude se compose, elle ne se résout pas.**
% Une transcription qui lisse un mot douteux a détruit la seule chose pour
% laquelle on la faisait — un lecteur doit pouvoir savoir, à chaque endroit,
% ce qui était sur le feuillet et ce qui a été ajouté. Les six macros qui
% suivent sont donc l'appareil critique, et non de la décoration.
%
% Les mêmes commandes sont reconnues par `scripts/render.mjs`, qui produit la
% vue de lecture du site. Une macro ajoutée ici doit l'être aussi là-bas,
% faute de quoi le rendu échouera bruyamment — ce qui est voulu : un rendu
% silencieusement faux est pire qu'un rendu absent.

\ProvidesPackage{grothendieck}[2026/08/08 Transcriptions du fonds Grothendieck]

\usepackage{amsmath,amssymb,amsthm}
\usepackage{mathrsfs}
\usepackage{stmaryrd}
% Les diagrammes commutatifs sont partout dans ce fonds : c'est la langue
% même de la géométrie algébrique de Grothendieck, et une transcription qui
% les rendrait en prose ne transcrirait rien.
\usepackage{tikz-cd}
% Français par défaut — c'est la langue des feuillets — mais l'anglais est
% chargé aussi : la traduction et le résumé sont des documents anglais, et la
% typographie française (espaces avant les deux-points) y serait une faute.
% Ces documents basculent par \selectlanguage{english} en tête de corps.
\usepackage[english,french]{babel}
\usepackage{fontspec}
\usepackage{geometry}
\geometry{a4paper,margin=25mm,marginparwidth=28mm}
\usepackage{marginnote}
\usepackage{xcolor}
% ulem, not soul. `soul`'s \st and \ul are built on a character-by-character
% scan that XeLaTeX's font machinery breaks: the document fails with an
% undefined control sequence rather than degrading. ulem does the same two jobs
% and is Unicode-engine safe. `normalem` stops it hijacking \emph, which the
% transcriptions use for Grothendieck's own emphasis.
\usepackage[normalem]{ulem}
\usepackage{hyperref}
\hypersetup{colorlinks=true,linkcolor=black,urlcolor=black,pdfborder={0 0 0}}

% --- Métadonnées du lot -----------------------------------------------------
% Reprises telles quelles par le script de rendu, qui les affiche en tête de
% la vue de lecture. Elles fixent ce que le fichier prétend couvrir : sans
% elles, une transcription flotte sans point d'attache dans un fonds de seize
% mille feuillets.

\newcommand{\folder}[1]{\gdef\grothFolder{#1}}
\newcommand{\batch}[1]{\gdef\grothBatch{#1}}
\newcommand{\leaves}[2]{\gdef\grothFirst{#1}\gdef\grothLast{#2}}
\newcommand{\foldertitle}[1]{\gdef\grothFoldertitle{#1}}
\gdef\grothFolder{?}\gdef\grothBatch{?}\gdef\grothFirst{?}\gdef\grothLast{?}\gdef\grothFoldertitle{}

% --- Le résumé --------------------------------------------------------------
%
% Il ouvre la lecture modernisée, sous le titre « Résumé », et il est composé
% en petit. Ce n'est pas un ornement : c'est le seul passage du document qui
% ne soit pas des mathématiques, et un lecteur doit pouvoir le reconnaître
% avant de l'avoir lu — puis le sauter s'il connaît déjà le sujet, ou s'y
% arrêter s'il ne le connaît pas. Le filet qui le suit marque la reprise du
% corps.

\newenvironment{resume}
  {\small\setlength{\parskip}{0.45em}}
  {\par\medskip\hrule\medskip}

% --- Mots-clés ---------------------------------------------------------------
%
% En anglais, en clôture du résumé de la lecture modernisée : le vocabulaire
% moderne sous lequel chercher ce que les feuillets construisent. Le manifeste
% du site (`npm run manifest`) les extrait pour étiqueter la cote entière —
% cette ligne est la source unique des tags, il n'y a pas de fichier à côté.
\newcommand{\keywords}[1]{\par\smallskip\noindent
  {\footnotesize\textsc{Keywords} --- \textit{#1}}\par}

% --- L'appareil critique ----------------------------------------------------

% Le numéro de feuillet, en marge. C'est l'ancre qui permet de régler un
% désaccord sur une formule en regardant *un* feuillet plutôt qu'une chemise
% de six cents. Le site s'en sert aussi pour tourner les pages du fac-similé
% quand on fait défiler la transcription.
\newcommand{\leaf}[1]{%
  \marginnote{\footnotesize\color{gray!70}\textsf{#1}}%
  \label{leaf:#1}\ignorespaces}

% Illisible : on ne devine pas. Un mot inventé qui se lit comme les autres est
% le pire résultat possible.
\newcommand{\ill}{\textcolor{red!60!black}{[\,\ldots\,]}}

% Lecture douteuse : le mot est donné, et signalé comme incertain.
\newcommand{\uncertain}[1]{\uline{#1}}

% Ajout de l'éditeur — une lettre manquante, un mot sous-entendu.
\newcommand{\add}[1]{\textcolor{blue!50!black}{[#1]}}

% Biffé par Grothendieck lui-même. Ce qu'il a rayé fait partie du document :
% une rature dit souvent où la pensée a changé de direction.
\newcommand{\struck}[1]{\sout{#1}}

% Note du transcripteur, sur l'état matériel du feuillet.
\newcommand{\note}[1]{\par\smallskip{\small\color{gray!60!black}\textit{[#1]}}\par\smallskip}

% Note marginale de Grothendieck, distincte du corps du texte.
\newcommand{\marginal}[1]{\par\smallskip\noindent\hspace{1em}%
  {\small\color{gray!60!black}#1}\par\smallskip}

% --- Titre ------------------------------------------------------------------

\AtBeginDocument{%
  \begin{center}
    {\large\bfseries Fonds Alexandre Grothendieck --- cote n\textsuperscript{o}~\grothFolder}\\[2pt]
    {\itshape\grothFoldertitle}\\[4pt]
    {\small feuillets \grothFirst--\grothLast\ (lot \grothBatch)}\\[2pt]
    {\footnotesize\color{gray!60!black}Transcription établie d'après le fac-similé numérisé
     par l'Université de Montpellier.\\
     Les lectures incertaines sont soulignées, les illisibles notées [\,\ldots\,],
     les ajouts entre crochets.}
  \end{center}
  \vspace{1em}}

\endinput


\folder{60}
\batch{1}
\pages{1}{20}
\dating{[à partir de 1962- à partir de 1971]}
\watermark{Édition de démonstration}
\foldertitle{Fonctions L / Équation fonctionnelle des fonctions L (cas géométrique) : notes manuscrites (s.d.)}

\begin{document}

\section{Notes sur l'équation fonctionnelle}
\note{titre écrit de sa main, souligné, en haut à droite de la chemise (page 1), qui ne porte rien d'autre}

\page{2}
Si $f$ est un endomorphisme d'un vectoriel $V$ de dimension finie $n$ sur le
corps $k$, on pose
\[
  P_f(t) = \det(1 - f t) = \prod_i (1 - \lambda_i t)
         = \sum_i (-1)^i \bigl(\operatorname{Tr} \lambda^i(f)\bigr)\, t^i
\]
les $\lambda_i$ étant les valeurs propres dans une \uncertain{clôture}
alg.\ close de $k$.

Donc les racines de $P_f(t)$ sont les inverses des valeurs propres non nulles
de $f$.

$P_f(t)$ peut (et doit) être considéré comme « additif » en $E$ (suite exacte
en $E$ \uncertain{transformée} en produits des polynômes). \ill{}
\[
  Z_f(t) = \frac{1}{P_f(t)} = \prod_i \frac{1}{1 - \lambda_i t}
         = \prod_i \sum_{\nu} \lambda_i^{\nu} t^{\nu}
         = \sum_{\nu} t^{\nu} \sum_{\nu_1 + \cdots + \nu_n = \nu}
             \lambda_1^{\nu_1} \cdots \lambda_n^{\nu_n}
         = \sum_{\nu} \operatorname{Tr}\bigl(\mathrm{Sym}^{\nu}(f)\bigr)\, t^{\nu} .
\]

D'autre part, $P_f(t)$ transforme le produit ($\otimes f$) en le « produit »
\ill{} \uncertain{formelles} \ill{} par $1$, si on a posé
\[
  (1 - a t) * (1 - b t) = 1 - a b t
\]
[et distributivité de $*$ \uncertain{p.r.}\ à $\cdot$, etc.\,\ldots]
\note{dans le coin inférieur droit, un petit schéma : trois rangées d'indices
($\rho-1$, $\rho$, $\rho+1$ en haut ; $-2$, $-1$, $0$ au milieu ; $\rho-1$,
$0$, $\rho+1$ et $\rho$, $\rho+1$, $\rho+2$ en bas), reliées par des ovales
obliques ; sa fonction n'est pas dite sur la page}

\page{4}
\[
  \left\lbrace
  \begin{array}{l}
    P_{f \otimes f'}(t) = P_f(t) * P_{f'}(t) \\[4pt]
    L_{f \otimes f'}(t) = \dfrac{1}{L_f(t) * L_{f'}(t)}
  \end{array}
  \right.
  \qquad
  L_{f_1 \otimes \cdots \otimes f_{\nu}}(t)
    = \Bigl(\prod L_{f_i}(t)\Bigr)^{(-1)^{\nu+1}}
\]
\note{le signe de produit de la dernière formule est tracé comme un $\prod$ ;
d'après les deux lignes de gauche il s'agit sans doute du produit $*$}

De même, pour « la » $\lambda$-structure de $1 + t\,\ill{}[[t]]^{+}$ où
\struck{$\lambda^p$} $\lambda^p$ transforme $\prod_i (1 - \lambda_i t)$ en
$\prod_{i_1 < \cdots < i_p} (1 - \lambda_{i_1} \cdots \lambda_{i_p} t)$, on
aura
\[
  \left\lbrace
  \begin{array}{l}
    P_{\lambda^i f}(t) = \lambda^i P_f(t) , \\[6pt]
    \dfrac{1}{L_{\lambda^i f}(t)} = \lambda^i \dfrac{1}{L_f(t)}
  \end{array}
  \right.
\]
donc, i.e.
\[
  -\sum (-1)^i \bigl[L_{\lambda^i f}(t)\bigr] s^i
    = \lambda^{*}\bigl(\bigl[-L_f(t)\bigr]\bigr)(-s)
    = \frac{1}{\lambda^{*}\bigl[L_f(t)\bigr](-s)}
    = \sum \sigma^i\bigl(L_f(t)\bigr) s^i
\]
\note{au-dessous, une ligne barrée de plusieurs traits :
$\bigl(\sum [\lambda^i(L_f(t))] s^i\bigr) * \bigl(\sum [L_{\lambda^i f}(t)] s^i\bigr) = -1 = \bigl[\tfrac{1}{1-t}\bigr]$ ;
un exposant $(-s)$ au dénominateur de la ligne précédente est aussi biffé.
Dans la première égalité, un facteur $(-1)^i$ est écrit puis surchargé ; la
lecture en reste douteuse}
\[
  \left\lbrace
  \begin{array}{l}
    L_{\lambda^i(f)}(t) = \bigl(\sigma^i L_f(t)\bigr)^{(-1)^{i+1}} \\[4pt]
    L_{\sigma^i(f)}(t) = \bigl(\lambda^i L_f(t)\bigr)^{(-1)^{i+1}}
  \end{array}
  \right.
\]

\page{6}
Considérons le cas d'un \emph{automorphisme} $f$, et voyons quel est le
rapport avec la contragrédiente $\check f$ :
\[
  P_{\check f}(t) = \prod (1 - \lambda_i^{-1} t)
  = \prod_i (-\lambda_i^{-1} t)(1 - \lambda_i t^{-1})
  = \frac{(-1)^n t^n}{\prod \lambda_i} \prod (1 - \lambda_i t^{-1})
  = \frac{(-1)^n t^n}{\det f}\, P_f(t^{-1})
\]
\note{la première égalité portait d'abord $\det(1 - \check f t)$, biffé ; une
ligne intermédiaire, biffée, commençait par une somme en $\lambda_i^{-1}$}
donc
\note{une ligne biffée : $P_{\check f}(t) = \frac{(-1)^n t^n}{\det f}
\ldots \varepsilon(f) \ldots P_f(t^{-1})$, et à droite « $\varepsilon(f) =$ »
laissé en suspens}
\[
  \boxed{P_{\check f}(t) = (-t)^{\varepsilon(f)} \, (\det f)^{-1} \, P_f(t^{-1})}
\]
i.e.
\[
  \boxed{L_{\check f}(t) = (-t)^{-\varepsilon(f)} \, \det f \; L_f(t^{-1})}
\]
\note{le signe $-$ de l'exposant $-\varepsilon(f)$ de la seconde formule
encadrée est douteux sur la page ; entre $\det f$ et $L_f$, un mot écrit au
crayon, repassé, ne se lit pas}

Si \ill{}, le deuxième membre \ill{} \ill{}, \ill{} ce que
\uncertain{signifierait} $P_f(t^{-1})$ pour une série formelle quelconque :
il est essentiel de se borner aux séries formelles qui sont
\uncertain{justement} obtenues à partir de \emph{polynômes} [à terme
constant $1$, à terme dominant inversible --- cette dernière condition n'étant
utile que si $k$ est un anneau général\ldots].

\page{8}
Si $E$ est un \struck{\ill{}} $\mathbb{Q}_\ell$-faisceau sur $\mathbb{F}_p$,
\struck{\ill{}} on trouve donc
\note{« $\mathbb{Q}_\ell$-faisceau » est écrit au-dessus du mot biffé, qui
pourrait être « motif »}
\[
  L_{\check E}(t) = (-t)^{-n} \det f_E \; L_E(t^{-1})
\]
où $n = \operatorname{rang}$ de $E$. Si $E$ est de poids $\rho$ et s'il s'agit
d'un \emph{motif} satisfaisant aux conj.\ de Weil
\note{« de poids $\rho$ » est ajouté dans une bulle au-dessus de la ligne}
\[
  \det f_E = \varepsilon(E)\, p^{\frac{n\rho}{2}}
\]
avec $\varepsilon = \pm 1$ [dépendant du motif $\overset{n}{\wedge} E$,
\struck{\ill{}} \uncertain{parfaitement} \ill{}] ; notons que $n\rho$ est
\ill{} \ill{} (i.e.\ $\rho$ impair $\Rightarrow$ $n$ pair).

On trouve donc
\[
  L_{\check E}(t) = \varepsilon(E) \Bigl(-\frac{p^{\rho/2}}{t}\Bigr)^{n} L_E(t^{-1}) .
\]

D'autre part, \uncertain{par dualité}, on aura :
\[
  \check E \simeq E(\rho)
\]
d'où
\[
  L_{\check E}(t) = \bigl(L_E(t) * L_{\mathbb{Q}_\ell(\rho)}(t)\bigr)^{-1}
  = \Bigl(L_E(t) * \frac{1}{1 - p^{-\rho} t}\Bigr)^{-1}
  = L_E(t) * (1 - p^{-\rho} t)
\]
\[
  L_{\check E}(t) = L_E\Bigl(\frac{t}{p^{\rho}}\Bigr)
\]

\page{10}
On trouve donc
\[
  \boxed{L_E\Bigl(\frac{1}{t p^{\rho}}\Bigr)
    = \det f_E \, (-t)^{n} \, L_E(t)}
\]
\note{le membre de droite commençait par un facteur biffé et surchargé,
élevé à la puissance $n$ ; l'argument de gauche portait d'abord un autre
symbole, surchargé en $\frac{1}{t p^{\rho}}$}
[\uncertain{rétablissant} le poids de $E$, qui s'introduit via
$\check E \simeq E(\rho)$]

Soit $E$ un $\mathbb{Q}_\ell$-faisceau constructible sur $X$
\struck{\ill{}} \add{propre} sur $\mathbb{F}_p$, alors on aura
\note{« propre » est écrit au-dessus du mot biffé}
\[
  \boxed{H^i_c(X, DE) \simeq D\bigl(H^{-i}(X, E)\bigr)}
\]
par le th.\ de dualité globale.
\struck{d'où dans
$\chi(DE) = \sum (-1)^i H^i(X, D(E)) \simeq D(\chi(E))$}
\note{ce passage est biffé de deux grandes boucles et d'un trait ; à la
place, une accolade renvoie à la formule suivante}
\[
  R f_{*}(DE) = D\bigl(R f_{*}(E)\bigr)
\]
d'où
\[
  (1)\quad
  \left\lbrace
  \begin{array}{l}
    \boxed{L_{DE}(t) = (-t)^{-\chi(E)}\, \delta(E)\, L_E(t^{-1})} \\[6pt]
    \chi(E) = \operatorname{rang} R f_{*} E = \sum (-1)^i \operatorname{rg} H^i(\bar X, \bar E) \\[6pt]
    \delta(E) = \prod_i \bigl(\det f_{H^i(\bar X, \bar F)}\bigr)^{(-1)^i}
      = \varepsilon(E)\, p^{\sum_i (-1)^i \frac{i + \rho}{2} b_i}
  \end{array}
  \right.
\]
\marginal{en face de $\chi(E)$ : invariant géométrique ; en face de
$\delta(E)$ : invariant arithmétique}
\marginal{sous l'égalité $\delta(E) = \varepsilon(E) p^{\ldots}$ : avec les
conjectures standard, si $X$ \struck{lisse} propre et lisse et $F$ c.t.\
\uncertain{fondamental} \uncertain{pur} de poids $\rho$}
\marginal{$b_i = \operatorname{rang} H^i(\bar X, \bar F)$}
\note{dans l'exposant de $p$, le numérateur $i + \rho$ est surchargé et la
lecture en reste douteuse ; la page écrit tantôt $E$, tantôt $F$ pour le
même faisceau}

\page{12}
\emph{Dans le cas favorable}
\[
  \sum_i (-1)^i \frac{i + \rho}{2} b_i
  = \frac{\rho}{2}\chi(E) + \frac{1}{2}\sum_{i=0}^{2n} (-1)^i i\, b_i
\]
\note{une accolade sous la seconde somme renvoie au calcul suivant}
\[
  \sum_{i=0}^{2n}(-1)^i i\, b_i
  = \sum_{i=0}^{n-1}\bigl[(-1)^i i\, b_i + (-1)^{2n-i}(2n-i)\, b_{2n-i}\bigr]
    + (-1)^n n\, b_n
\]
comme $b_i = b_{2n-i}$ :
\[
  = n\sum_{i=0}^{n-1} (-1)^i b_i + (-1)^n n\, b_n
  = n \sum_{i=0}^{2n} (-1)^i b_i = n\,\chi(E)
\]
\note{la page écrit $2n\sum_{i=0}^{n-1}$ au début de la dernière ligne, et
un facteur surchargé devant $\sum_{i=0}^{2n}$ ; on donne la ligne telle
qu'elle aboutit, $n\,\chi(E)$}
\[
  \sum_i (-1)^i \frac{i + \rho}{2} b_i = \frac{n + \rho}{2}\,\chi(E)
\]
donc
\note{un premier membre de droite est biffé avant $\frac{n+\rho}{2}\chi(E)$}
\[
  \delta(E) = \varepsilon(E)\, p^{\frac{n+\rho}{2}\chi(E)} ,
  \qquad \varepsilon(E) = \pm 1
\]
Par suite l'équation fonctionnelle devient dans ce cas
\[
  (2)\quad
  L_{DE}(t) = \varepsilon(E)\Bigl(-\frac{p^{\frac{n+\rho}{2}}}{t}\Bigr)^{-\chi(E)} L_E(t^{-1})
\]
\note{le numéro (2) surcharge un autre chiffre ; le $t$ du dénominateur
surcharge un $1$}

On remarque, si on admet que
\[
  \check E \simeq E(\rho)
\]
globalement sur $X$ \struck{\ill{}}
[il suffit que ce soit vrai sur le \uncertain{complémentaire} d'une
\ill{} de $X$, \ill{} \ill{} \ill{}]
\[
  D(E) \simeq \check E(n)[2n] \simeq E(\rho + n)[2n]
\]
\note{sous les deux crochets $[2n]$, deux flèches légendées « shift » et
« shift » ; avant $D(E)$ un symbole biffé ; entre $\simeq$ et
$\check E(n)$, un $\check E \otimes \mathbb{Q}_\ell(n)$ biffé}
et par suite \struck{$L_{DE} \simeq L_{E(\rho+n)} = L_{Rf_{*}E(\rho+n)}$}
\[
  R f_{*}(D(E)) \simeq R f_{*}(E)(\rho + n)[2n]
\]
\note{dans $(\rho + n)$ le $n$ surcharge un autre signe}
d'où
\[
  L_{DE}(t) = L_{Rf_{*}(D(E))}(t) \simeq L_{Rf_{*}(E)}\bigl(p^{-(\rho+n)} t\bigr)
\]
\struck{et il \ill{}}

\page{14}
\[
  L_{DE}(t) = L_E\bigl(p^{-(n+\rho)} t\bigr)
\]
d'où, portant dans (2) et remplaçant $t$ par $t^{-1}$ :
\[
  (3)\quad
  \boxed{L_E\Bigl(\frac{1}{p^{n+\rho}\, t}\Bigr)
    = \varepsilon(E)\bigl(-p^{\frac{n+\rho}{2}} t\bigr)^{\chi(E)} L_E(t)}
\]
Posant
\[
  \boxed{\xi_E(t) = t^{\chi(E)/2} L_E(t)}
\]
il vient l'équation fonctionnelle
\[
  \boxed{\xi\Bigl(\frac{1}{p^{n+\rho}\, t}\Bigr) = \varepsilon'(E)\, \xi(t)}
\]
\[
  \varepsilon'(E) = \pm 1 , \qquad
  \varepsilon'(E) = (-1)^{\chi(E)} \prod_i \operatorname{signe} \det f_{H^i(\bar X, \bar F)}
  = (-1)^{\chi(E)} \operatorname{signe} \det f_{H^{*}(\bar X, \bar F)}
\]
\struck{\ill{} \ill{} \ill{} \ill{} l'équation fonctionnelle \ill{}
\ill{} $\varepsilon(E)$}

Notons d'ailleurs que $H^i(\bar X, \bar E)$ et
$H^{2n-i}(\bar X, \bar{\check E}) \simeq H^{2n-i}(\bar X, \bar E(\rho))
= H^{2n-i}(\bar X, \bar E)(\rho)$ [étant \uncertain{dual} l'un de l'autre
\uncertain{à valeurs dans} « $\mathbb{Q}_\ell(n)$ »], et l'opération de $f$
sur \add{$\mathbb{Q}_\ell(n)$} \uncertain{étant} \ill{} \ill{}
\uncertain{positive}, $\varepsilon_i(E) = \varepsilon_{2n-i}(E)$, et par suite
on trouve que
\[
  \boxed{\varepsilon(E) = \operatorname{signe} \det f_{H^n(\bar X, \bar E)}} .
\]
\note{deux lignes biffées précèdent « Notons » ; « $\mathbb{Q}_\ell(n)$ » est
ajouté au-dessus de la ligne}
Par la même raison \uncertain{analogue}, $(-1)^{\chi(E)} = (-1)^{b_n}$.
Donc
\[
  \varepsilon'(E) = (-1)^{b_n(E)} \operatorname{signe} \det f_{H^n(X, E)}
\]

\page{16}
\emph{Si $n + \rho$ est impair}, alors $b_n = \dim H^n(\bar X, \bar E)$ est
pair, donc $(-1)^{b_n} = 1$. De plus, un raisonnement bien connu de
\uncertain{Serre} \ill{} \ill{} \uncertain{(forme alternée non dégénérée
invariante)} montre que dans ce cas $\operatorname{signe} \det f_{H^n(\bar X, \bar E)}$
est $+1$. Donc on trouve alors $\varepsilon'(E) = 1$.
\note{sous $H^n(\bar X, \bar E)$, une accolade et deux mots, dont le second
est $n$ ou $\rho$, qui ne se lisent pas}

Si $n + \rho$ est pair, on ne peut déterminer \uncertain{(\ill{} de $2$
\ill{})} le signe de façon absolue, mais on peut seulement dire que si
$\varepsilon_1$ ($\varepsilon_{-1}$) est la \add{multiplicité}
\struck{\ill{}} de \struck{\ill{}} la valeur propre $1$ ($-1$) de $f$ sur
$H^n(\bar X, \bar F)$, on a
\note{« multiplicité » est écrit au-dessus du mot biffé}
\[
  (-1)^{b_n} = (-1)^{\varepsilon_1 + \varepsilon_{-1}} , \qquad
  \varepsilon(E) = (-1)^{\varepsilon_{-1}} ,
\]
donc
\[
  \boxed{\varepsilon'(E) = (-1)^{\varepsilon_1}} .
\]

Lorsque l'on n'est pas \uncertain{dans} le cas favorable ($X$ lisse, $E$
\ill{}, $\check E \simeq E(\rho)$ \emph{globalement}), \uncertain{donc lorsque}
$DE$ ne s'exprime plus de façon aussi simple en fonction de $E$, l'équation
fonctionnelle (1) ne s'exprime plus \uncertain{uniquement} en termes de la
seule fonction $L_E$ [mais relie \emph{deux} fonctions $L_{DE}$, $L_E$].
\note{avant $\check E$ un symbole surchargé}

\page{18}
$X$ lisse purement de dim.\ $n$,

$E$ et tordu purement de poids $\rho$,

$X \xrightarrow{f} \operatorname{Spec} \mathbb{F}_p$ compactifiable.
\note{« $\operatorname{Spec}$ » est surchargé}

\emph{Dualité}
\[
  D\, R f_{!}(E) = R f_{*}(DE)
\]
Or par hypothèse sur $E$
\[
  DE = \check E(n)[2n] \simeq E(\rho + n)[2n]
\]
\note{dans $E(\rho + n)$, le $n$ surcharge une autre lettre}
D'où
\[
  D\, R f_{!}(E) \simeq \bigl(R f_{*}(E)\bigr)(\rho + n)[2n]
\]
d'où en prenant les classes, et écrivant
\[
  \begin{aligned}
    \operatorname{cl} D\bigl(R f_{!}(E)\bigr)
      &= \bigl(\operatorname{cl} R f_{*}(E)\bigr)(\rho + n) \\
      &= \operatorname{cl}\bigl(R f_{!}(E)\bigr)(\rho + n)
         + \operatorname{cl} R f_{\infty}(E)(\rho + n)
  \end{aligned}
\]
\note{à la première ligne, un exposant est biffé devant $(\rho + n)$ ; la
lecture de l'indice $\infty$ de $R f_{\infty}$ est celle de la page}
d'où, pour les fonctions $L$
\struck{$L_{DE}$}
\[
  (-t)^{-\chi(E)}\, \delta(E)\, L_E\Bigl(\frac{1}{t}\Bigr)
  = L_E\bigl(p^{-(\rho+n)} t\bigr)\, L^{\infty}_E\bigl(p^{-(\rho+n)} t\bigr)
\]
\note{au début de la ligne, le $L_{DE}$ est biffé et surchargé ; un autre
symbole biffé précède $\delta(E)$ ; l'exposant de $L^{\infty}_E$ est
surchargé}
ou encore
\note{la page s'arrête sur ces mots}

\page{20}
\[
  \boxed{L_E\Bigl(\frac{1}{p^{n+\rho}\, t}\Bigr)
    = A(t)\, \delta(E)\, (-t)^{\chi(E)} L_E(t)}
\]
\marginal{dans un cadre, en haut à droite : en admettant la résolution des
singularités \ill{} \ill{} et les conjectures de Weil \ldots}
\[
  \left\lbrace
  \begin{array}{l}
    A(t) = L^{\infty}_E\Bigl(\dfrac{1}{p^{n+\rho}\, t}\Bigr)^{-1} \\[8pt]
    \chi(E) = \sum (-1)^i \operatorname{rang} H^i_!(\bar X, \bar F) \\[6pt]
    \delta(E) = \prod_i \bigl(\det f_{H^i_!(\bar X, \bar F)}\bigr)^{(-1)^i}
  \end{array}
  \right.
\]
\note{dans $A(t)$, un $p$ biffé précède la fraction ; sous $A(t)$ un signe
$\simeq$ reste sans suite ; dans $\delta(E)$, un mot biffé précède
« $\det$ »}

\emph{NB} $L^{\infty}_E(t)$ ($= L_{R f_{\infty}(E)}(t) =
L_{R i_{*}(E)|Y}(t)$) est de \add{type} \struck{\ill{}} $\leq \tau + (n-1)$
($\tau$ étant le type de $E$) et de poids $\leq \rho + 2n$.
\note{« type » est écrit au-dessus du mot biffé ; au premier membre, $\tau$
surcharge une autre lettre}

\emph{Si $X$ complet}, $A(t) = 1$.
\note{un long trait oblique sépare ce bloc de la suite}

On peut aussi écrire
\[
  \chi(E) = \varepsilon\bigl(R f_{!}(E)\bigr) , \qquad
  \delta(E) = \delta\bigl(R f_{!}(E)\bigr)
\]
\note{dans la seconde égalité, un « $\det$ » est biffé devant $\delta$}
\[
  \varepsilon\bigl(R f_{!}(E)\bigr) = \varepsilon\bigl(D R f_{!}(E)\bigr) ,\qquad
  \delta\bigl(R f_{!}(E)\bigr) = \delta\bigl(D(R f_{!}(E))\bigr)^{-1}
\]
\[
  D R f_{!}(E) = R f_{*}(DE) = \bigl(R f_{*}(E)\bigr)(\rho + n)[2n]
\]
\note{la page note $\varepsilon$ dans la première égalité là où l'on
attendrait $\chi$ ; on transcrit tel quel. Devant $\det$ et devant
$(\rho + n)$ des signes sont surchargés ; « d'où d'où », le second biffé,
mène à la suite}
d'où
\[
  \left\lbrace
  \begin{array}{l}
    \chi(E) = \sum (-1)^i \operatorname{rang} R^i f_{*}(E) \\[6pt]
    \delta(E) = p^{+(\rho+n)\chi(E)}\, \delta\bigl(R f_{*}(E)\bigr)^{-1}
      = p^{(n+\rho)\chi(E)} \prod_i \bigl(\det f_{H^i(\bar X, \bar F)}\bigr)^{(-1)^{i+1}}
  \end{array}
  \right.
\]
\note{l'accolade est précédée d'un signe barré}
[Comparant, lorsque $X$ est complet, les deux \uncertain{valeurs} obtenues de
$\delta(E)$, on trouve dans ce cas
\[
  \delta(E)^2 = p^{(n+\rho)\chi(E)}
\]
($X$ complet)
ce qui est bien compatible avec les conjectures de Weil sur les
val.\ absolues des valeurs propres\ldots]


\end{document}
